\documentclass[10pt,a4paper]{amsart}
\usepackage[margin=19mm]{geometry}
\usepackage{amsmath,amssymb,mathtools}
\usepackage[scr=rsfs,cal=euler]{mathalfa}
\usepackage{xfrac}
\usepackage[unicode,hidelinks,bookmarks=false]{hyperref}
\newcommand{\ie}{i.e.\ }
\newcommand{\cf}{cf.\ }
\newcommand{\resp}{respectively\ }
\newcommand{\Subsection}{\subsection}
\providecommand{\nobreakdash}{\nobreak\hbox{-}}
\setlength{\emergencystretch}{2em}
\multlinegap0pt
\usepackage{bm}
\usepackage{bbm}
\usepackage{dsfont}
\usepackage{xspace}
\usepackage{cancel}
\usepackage{mathtools}
\usepackage{amsthm}
\makeatletter
\@ifundefined{theorem}{\newtheorem{theorem}{Theorem}}{}
\@ifundefined{corollary}{\newtheorem{corollary}[theorem]{Corollary}}{}
\@ifundefined{lemma}{\newtheorem{lemma}[theorem]{Lemma}}{}
\makeatother
\title[Function-Correcting Codes for Insertion-Deletion Channel]{Function-Correcting Codes for Insertion-Deletion Channel}
\author{Anamika Singh, Abhay Kumar Singh}
\date{Source: arXiv:2512.07243v3 (CC BY 4.0)}
\begin{document}
\maketitle

\noindent\emph{Source and licence.} Function-Correcting Codes for Insertion-Deletion Channel. Version arXiv:2512.07243v3 (CC BY 4.0), distributed under the Creative Commons Attribution 4.0 International licence, as recorded in the arXiv licence metadata for this exact version and shown on the abstract page https://arxiv.org/abs/2512.07243v3. Full rights notice: licenses/arxiv-2512.07243-CC-BY-4.0.txt.\par\smallskip

\noindent\textit{Editorial scope.} Counted unit: the Lemma whose source label is \texttt{max\_lower\_bound} in the article (source0.tex lines 2030--2048, source label \texttt{max\_lower\_bound}), delivered together with the article's own complete original proof of it (source0.tex lines 2049--2191, one \texttt{proof} environment of 143 lines). No mathematics is written, repaired or re-derived: statement and proof are the article's own text line for line. Required context retained uncounted: run\_insdel\_bound (lines 330--334); lower\_bound\_red (lines 1052--1057). Every definition, display and label of those lines is retained unchanged. Omitted from this package and not redistributed: 1--329, 335--1051, 1058--2029, 2192--2578. Nothing inside a retained excerpt is omitted or altered: every retained body is delivered exactly as the article's lines are written, so no omission receipt of any kind accompanies this package. Citations: the retained excerpts carry no citation command at all, so no bibliography entry of the article is transcribed and no reference is redistributed. Reference adaptation: none was needed and none was made; every reference inside the delivered unit resolves inside it, so no unresolved reference or citation is produced. Printed slips kept literal: no printed word, symbol, hypothesis or number of the counted statement or of its proof was altered. Layout and class: the article's own class is \texttt{article} with packages including threeparttable, amsmath, amssymb, geometry, amsthm, url, tikz, inputenc, babel, graphicx, hyperref, adjustbox, makecell, bm; the wrapper is the lane's pinned \texttt{amsart} editorial wrapper at 10pt on A4, which restates the theorem-like environments the retained text needs and adds the article's own macro definitions that the retained text needs (none), with extra wrapper packages (bm, bbm, dsfont, xspace, cancel, mathtools). The delivered numbering is the unit's own. Assets: no retained span contains a figure environment, a graphics-inclusion command, a PDF-page inclusion, a table or a TikZ picture; no image file is copied into this package, the package declares no asset dependency and no image file is redistributed with it. Licence: version arXiv:2512.07243v3 of this article is distributed under the Creative Commons Attribution 4.0 International licence, as recorded in the arXiv licence metadata for this exact version and shown on the abstract page https://arxiv.org/abs/2512.07243v3. A full rights notice is delivered at licenses/arxiv-2512.07243-CC-BY-4.0.txt. Characters: every retained excerpt is pure ASCII, so the delivered engine, XeLaTeX with its default Latin Modern text font, reports no missing character.
    \begin{lemma}\label{run_insdel_bound}
        Let $x,y \in \mathbb{F}_2^n$ be two binary vectors of length $n$ and let $r(x)$ and $r(y)$ be the number-of-runs corresponding to vectors $x$ and $y$ respectively. Then,
        $$ d_{ID}(x,y) \ge 2 \Big\lceil \frac{|r(y) - r(x)|}{2} \Big\rceil.$$
        Moreover, when $|r(x) - r(y)|$ is odd then $d_{ID}(x,y)\geq |r(x) - r(y)| + 1.$
    \end{lemma}

    \begin{corollary}\label{lower_bound_red}
        Consider a collection of $M$ distinct binary words $x_1, x_2, \ldots, x_M$ of length $k$. Then, for a function $f$, the optimal redundancy of an FCIDC is lower bounded as follows:
        $$ r^f_{ID}(k, t) \geq N^{(1)}_{ID}(\boldsymbol{I}_f^{(1)}(t, x_1,x_2, \ldots,x_{M})).$$
        For a non-constant function $f$, 
        $$ r^f_{ID}(k, t) \geq N^{(1)}_{ID}(2, 2t) = t. $$
    \end{corollary}

    \begin{lemma}\label{max_lower_bound}
        Let $f(x) = r_{max}(x)$ and $x_i =0^i(10)^{(k-i)/2}$ if $k-i$ is even and $x_i =0^i(10)^{(k-i-1)/2}1$ if $k-i$ is odd for $i \in \{1,2,\ldots,k\}$.Then,
        
          \[
            r_{ID}^f(k,t) \geq N_{ID}^{(1)}(\boldsymbol{I}_f^{(1)}(t, x_1, x_2, \ldots, x_k)).
          \]  
        
        where
        
        \[
            [\boldsymbol{I}_f^{(1)}(t, x_1, x_2, \ldots, x_k)]_{ij}=
            \begin{cases}
                0,  & i = j, \\[4pt]
                2t + 1 - |i - j| , & i \neq j \text{ and } |i-j| \text{ is odd },\\
                2t + 2 - |i-j|,    & i \neq j \text{ and } |i-j| \text{ is even }. 
            \end{cases}
        \]
        
    \end{lemma}

    \begin{proof}
        Consider the following set of representative vectors: $x_i =0^i(10)^{(k-i)/2}$ if $k-i$ is even and $x_i =0^i(10)^{(k-i-1)/2}1$ if $k-i$ is odd for $i \in \{1,2,\ldots,k\}$. Then, $f(x_i) = r_{max}(x_i) = i$. We claim that the insdel distance between these representative set of vectors is given by:
        \[
            d_{ID}(x_i, x_j) =
            \begin{cases}
                |i-j|, & \text{ if $|i-j|$ is even}\\
                |i - j| + 1, & \text{if $|i-j|$ is odd}
            \end{cases}
        \]
        WLOG assume $k$ is even and $i > j$.\\
        \noindent \textbf{Case 1:} Both $i$ and $j$ are even.
        
        \smallskip
        \noindent Then, $LCS(x_i,x_j) \geq j + \frac{i-j}{2} + k- i = k - \frac{i-j}{2}.$\\
        Hence, $d_{ID}(x_i,x_j) = 2(k - LCS(x_i,x_j)) \leq i-j$.\\
        From Lemma~\ref{run_insdel_bound} we have $d_{ID}(x_i,x_j) \geq |r(x_i) - r(x_j)| = |(1 + k -i) -(1 + k - j)| = |j-i| = i-j.$\\
        Hence, $d_{ID}(x_i, x_j) = i-j$.\\
        
        \noindent \textbf{Case 2:} Both $i$ and $j$ are odd.
        
        \smallskip
        \noindent Then, $LCS(x_i,x_j) \geq j + \frac{i-j}{2} + ( k- i-1) +1 = k - \frac{i-j}{2}.$\\
        Therefore, using the same argument as in Case 1, we get $d_{ID}(x_i, x_j) = i-j$.\\
        
        % \noindent \textbf{Case 3:} $i$ is even $j$ is odd.\\
        % Then, $LCS(x_i,x_j) \geq j + \frac{i-j-1}{2} -1+ ( k- i-1) = k - \frac{i-j-1}{2} + 1.$\\
        % Hence, $d_{ID}(x_i,x_j) = 2(k - LCS(x_i,x_j)) \leq i-j + 1$.\\
        % From Lemma~\ref{run_insdel_bound} we have $d_{ID}(x_i,x_j) \geq |r(x_i) - r(x_j)| + 1= |(1 + k -i) -(1 + k - j)| + 1= |j-i| +1 = i-j + 1.$\\
        % Therefore, using the same argument as in Case 1, we get $d_{ID}(x_i, x_j) = i-j+1$.\\
        
        % \noindent \textbf{Case 4:} $i$ is odd $j$ is even.\\
        % Using the same argument of Case 3, one can easily verify that $d_{ID}(x_i, x_j) = i-j+1$.\\
        % Therefore, we can conclude that
        
        
        % Here is a compact, academically toned version of **Case 4**, written to mirror the structure of **Case 3** while keeping it self-contained and concise.

% ---

% \ 
% (so $i - j$ is odd, with $i - j \geq 1$).

% \smallskip
% {Lower bound on $d_{ID}(x_i, x_j)$.}  
% From the definitions,  
% $x_i = 0^{i}(10)^{(k-i-1)/2}1$ and $x_j = 0^{j}(10)^{(k-j)/2}$, giving  
% \[
% r(x_i) = 1 + (k - i), \qquad r(x_j) = 1 + (k - j),
% \]  
% so $|r(x_i) - r(x_j)| = i - j$, which is odd. By Lemma~\ref{lem:runs},  
% \[
% d_{ID}(x_i, x_j) \;\geq\; i - j + 1.
% \]

% \smallskip
% {Upper bound on $d_{ID}(x_i, x_j)$.}  
    \noindent\textbf{Case 3.} $i$ is odd and $j$ is even.\\
    Then, $x_i = 0^i\,(10)^{(\frac{k-i-1}{2})}1$ and $x_j = 0^j\,(10)^{(\frac{k-j}{2})}$.
    Consider a vector $z \;=\; 0^{(i+j-1)/2}\,(10)^{(k-i-1)/2}1 \in \mathbb{F}_2^{k - \frac{i-j+1}{2}}$. We claim that $z$ is a common subsequence of both $x_j$ and $x_i$.
    
    \medskip
    \begin{itemize}
    \item {$z$ is a subsequence of $x_j$:}      
     $z$ can be obtained from $x_j$ by deleting $\frac{i-j+1}{2}$ bits from $x_j$. Delete the first $(i - j - 1)/2$ ones of $x_j$ at positions $j+1, j+3, \ldots, i-2$ and then delete the last zero of $x_j$ at position $k$. After these deletions, $x_j$'s remaining characters at positions $\{1, \ldots, j\} \cup \{j+2, j+4, \ldots, i-1\} \cup \{i, i+1, \ldots, k-1\}$ are
    \[
    \underbrace{0, \ldots, 0}_{j}\;\underbrace{0, \ldots, 0}_{(i-j-1)/2}\;\underbrace{(10)^{(k-i-1)/2}\,1}_{\text{positions } i, \ldots, k-1}
    \;=\; 0^{(i+j-1)/2}\,(10)^{(k-i-1)/2}\,1 \;=\; z.
    \] 

    
    \item {$z$ is a subsequence of $x_i$:}  
    Since $j \leq i - 1$, we have $(i+j-1)/2 \leq i - 1 < i$, so $x_i$'s leading block $0^{i}$ contains strictly more than $(i+j-1)/2$ zeros. Match $z$'s leading $(i+j-1)/2$ zeros with $x_i$'s 0 at positions $1, 2, \ldots, (i+j-1)/2$, and match $z$'s tail $(10)^{(k-i-1)/2}\,1$ with $x_i$'s tail $(10)^{(k-i-1)/2}\,1$ at $x_i$ positions $i+1, i+2, \ldots, k$. The matched positions form the strictly increasing sequence $1, \ldots, (i+j-1)/2, i+1, \ldots, k$ because $(i+j-1)/2 < i + 1$, so $z$ is a subsequence of $x_i$.
    \end{itemize}
    Hence $\mathrm{LCS}(x_i, x_j) \geq k - \frac{i-j+1}{2}$, yielding  
    \[
    d_{ID}(x_i, x_j)  \leq i - j + 1.
    \]  
    By Lemma~\ref{run_insdel_bound}, $d_{ID}(x_i,x_j)\ge |r(x_i)-r(x_j)| = |(1+k-i)-(1+k-j)| = i-j$. Since $(x_i,x_j)$ have equal length, $d_{ID}(x_i,x_j)$ is even, while $i-j$ is odd, we have $d_{ID}(x_i,x_j)\ge i-j+1$. Combined with the upper bound, $d_{ID}(x_i,x_j)=i-j+1$. 

    \smallskip

    \noindent \textbf{Case 4:} $i$ is even and $j$ is odd.\\
    
    Then
    \[
    x_i = 0^{i}\,(10)^{\frac{k-i}{2}}, \qquad 
    x_j = 0^{j}\,(10)^{\frac{k-j-1}{2}}\,1.
    \]

    
    Consider the word
    \[
    z = 0^{\frac{i+j-1}{2}}\,(01)^{\frac{k-i}{2}} \;\in\; \mathbb{F}_2^{\,k - \frac{i-j+1}{2}}.
    \]
    \begin{itemize}
    \item {$z$ is a subsequence of $x_i$:} 
    Match the leading $\frac{i+j-1}{2}$ zeros of $z$ with the first $\frac{i+j-1}{2}$ zeros of $x_i$ 
    (positions $1,\dots,\frac{i+j-1}{2}$). 
    For the trailing block $(01)^{\frac{k-i}{2}}$, match the $t$-th ``0'' with the zero at position 
    $i+2(t-1)$ of $x_i$ and the following ``1'' with the one at position $i+2t-1$ of $x_i$ 
    ($t=1,\dots,\frac{k-i}{2}$). 
    The matched indices are strictly increasing because $\frac{i+j-1}{2} < i$; hence $z$ is a subsequence of $x_i$.

    \item {$z$ is a subsequence of $x_j$:} 
    Delete from $x_j$ the first $\frac{i-j+1}{2}$ ones (positions $j+1, j+3, \dots, i$). 
    The remaining symbols consist of the zeros of the leading block (positions $1,\dots,j$), 
    the zeros at positions $j+2, j+4, \dots, i-1$ (there are $\frac{i-j-1}{2}$ of them), 
    and the trailing block $(01)^{\frac{k-i}{2}}$ (which starts at position $i+1$ because $i+1$ is odd and is a zero). 
    Concatenating these yields $z$. Thus $z$ is a common subsequence.
    \end{itemize}
    Therefore $\mathrm{LCS}(x_i,x_j) \ge |z| = k - \frac{i-j+1}{2}$, and
    \[
    d_{ID}(x_i,x_j) \le 2\bigl(k - \mathrm{LCS}(x_i,x_j)\bigr) \le i - j + 1.
    \]


    The run counts are $r(x_i)=1+(k-i)$, $r(x_j)=1+(k-j)$, so $|r(x_i)-r(x_j)| = i-j$, which is odd. 
    By Lemma~\ref{run_insdel_bound},
    \[
    d_{ID}(x_i,x_j) \ge i - j + 1.
    \]

    \noindent
    Combining the upper and lower bounds gives $d_{ID}(x_i,x_j) = i - j + 1$.


    \noindent
    Thus, in both Cases 3 and 4, when $i$ and $j$ have opposite parity,  
    \[
        d_{ID}(x_i, x_j) = |i - j| + 1.
    \]
    \[
            [\boldsymbol{I}_f^{(1)}(t, x_1, x_2, \ldots, x_k)]_{ij}=
            \begin{cases}
                0,  & i = j, \\[4pt]
                [2t + 1 - |i - j|]^+ , & i \neq j \text{ and } |i-j| \text{ is odd },\\
                [2t + 2 - |i-j|]^+,    & i \neq j \text{ and } |i-j| \text{ is even }. 
            \end{cases}
        \]
        
        and the lower bound on the redundancy, $r_{ID}^f(k,t) \geq N_{ID}^{(1)}(\boldsymbol{I}_f^{(1)}(t, x_1, x_2, \ldots, x_k))$ follows from Corollary~\ref{lower_bound_red}.
    
    \end{proof}

\end{document}
